\chapter{Configuration Reference}
\label{ch:configuration}

Blackship reads TOML from \file{blackship.toml} by default. Select another file
with the global \command{--config} option. Armada can merge several manifests in
order; later files override earlier values.

\section{Global configuration}

\begin{lstlisting}[language=TOML]
[config]
project = "orbital"
data_dir = "/var/blackship"
releases_dir = "/var/blackship/releases"
cache_dir = "/var/cache/blackship"
mirror_url = "https://download.freebsd.org/releases"
zfs_enabled = true
zpool = "zroot"
dataset = "blackship"
bootstrap_archives = ["base", "lib32"]
usage_sample_secs = 30
monitor_ping_url = "https://monitor.example/ping/token"
\end{lstlisting}

\begin{shiptable}{L{0.23\textwidth}L{0.18\textwidth}L{0.49\textwidth}}
\toprule
Key & Typical value & Meaning \\
\midrule
\setting{project} & string & Prefix for effective jail names. If omitted,
Blackship uses the manifest directory name, then the working-directory name,
with a path-derived Black Ship name as the final fallback. \\
\setting{data\_dir} & path & Persistent state, roots, scopes, audit, networks,
leases, and forward records. \\
\setting{releases\_dir} & path & Bootstrapped and committed release roots. \\
\setting{cache\_dir} & path & Downloads and reusable build material. \\
\setting{mirror\_url} & URL & Base URL for FreeBSD release acquisition. \\
\bottomrule
\end{shiptable}

\begin{shiptable}{L{0.23\textwidth}L{0.18\textwidth}L{0.49\textwidth}}
\multicolumn{3}{l}{\textit{Global configuration (continued)}} \\[0.35em]
\toprule
Key & Typical value & Meaning \\
\midrule
\setting{zfs\_enabled} & boolean & Enable dataset-backed provisioning and ZFS features. \\
\setting{zpool} & string & Pool containing the Blackship dataset tree. \\
\setting{dataset} & string & Dataset path below the pool. \\
\setting{bootstrap\_archives} & array & Release components to acquire; the
default is \setting{["base"]}. \\
\setting{usage\_sample\_secs} & integer & Supervisor racct and telemetry sample
period; zero disables sampling. \\
\setting{monitor\_ping\_url} & URL & Optional dead-man ping used by the supervisor. \\
\bottomrule
\end{shiptable}

When ZFS is enabled, the jails dataset is mounted at the configured data
directory's \file{jails} path. This keeps effective root paths consistent
between directory and dataset deployments.

Operational defaults are tunable in nested tables. The following example also
shows bridge VLAN filtering, which requires FreeBSD 15.0 or newer:

\begin{lstlisting}[language=TOML]
[config.rate_limit]
jail_start_capacity = 3.0
health_capacity = 5.0
health_refill_rate = 0.5

[config.health]
interval = 30
timeout = 10
start_period = 60
retries = 3

[config.retry]
base_delay_ms = 1000
max_delay_ms = 30000
multiplier = 2.0
max_attempts = 5
jitter_factor = 0.25

[config.bridge]
name = "vswitch1"
vlan_filtering = true

[config.bridge.trunk]
interface = "igb1"
tagged = [10, 20, 99]
disable_hwfilter = false
\end{lstlisting}

Rate limits bound concurrent starts and health-check dispatch. The
\setting{config.health} schema records the intended global check defaults; in
version 0.3.0, checks use their own values or the same built-in defaults shown
above. Retry values control exponential backoff for HTTP release operations;
jitter is expressed as a fraction from zero through one. The optional trunk
block declares the physical interface and VLANs admitted by a filtering bridge.

\section{Jail definitions}

Each service is an array entry:

\begin{lstlisting}[language=TOML]
[[jails]]
name = "api"
release = "15.1-RELEASE"
hostname = "api.internal"
depends_on = ["database", "cache"]
tags = ["web", "production"]
devfs_ruleset = 4
init = true
stop = "service api stop"
ephemeral = false
\end{lstlisting}

\setting{name} is required and must be safe for a FreeBSD jail name.
\setting{release} selects the seed userland. \setting{path} may replace the
default root path, and \setting{build} or \setting{jailfile} can point Armada at
a build context. \setting{hostname} defaults from the service name.

\setting{depends\_on} defines graph edges. \setting{tags} groups several jails
under a selector. \setting{init = true} runs the FreeBSD rc lifecycle. A
manifest \setting{stop} command overrides the Jailfile's \command{STOP}.

Additional jail parameters are accepted in \setting{params}, but identity and
isolation keys managed by Blackship---including name, path, VNET, securelevel,
devfs, address, and related \command{allow.*} controls---cannot be overridden
through an escape hatch. Use first-class configuration where provided.

\section{Resource controls}

\begin{lstlisting}[language=TOML]
[jails.resources]
memory = "2g"
vmemory = "4g"
cpu = "200"
maxproc = 512
maxthreads = 1024
openfiles = 4096
swap = "1g"
cores = 2
# cpuset = "0,1"
# rules = ["readbps:deny=50m/s"]
\end{lstlisting}

Named fields become rctl rules. \setting{cpu} is percentage-style pcpu policy;
values above 100 permit more than one core of aggregate CPU. \setting{cores}
asks Blackship to choose the least-loaded CPUs while \setting{cpuset} pins an
explicit list and wins when both are present. The \setting{rules} array accepts
advanced raw RCTL resource/action expressions; a jail subject is added when
omitted, and a rule targeting a different jail is rejected.

RCTL amount fields accept integers or binary size suffixes
\setting{k}, \setting{m}, \setting{g}, \setting{t}, \setting{p}, and
\setting{e}, optionally followed by \setting{b}. Blackship expands these to
the integer form required by the kernel before issuing the rule syscall.

If configured resource enforcement cannot be applied, startup fails and rolls
back. This avoids silently running a workload beyond its declared boundary.

\section{Network configuration}

\begin{lstlisting}[language=TOML]
[jails.network]
vnet = true
networks = ["frontend"]
ip = "10.20.0.15"
# ip_cidr = "10.20.0.15/24" # alternative with explicit prefix
gateway = "10.20.0.1"
mac_address = "58:9c:fc:00:20:15"
vlan_id = 20
backend = "epair"

[jails.network.dns]
mode = "custom"
nameservers = ["10.20.0.53", "1.1.1.1"]
search = ["internal.example"]
domain = "internal.example"
\end{lstlisting}

A named network supplies subnet, bridge, gateway, and backend runtime data.
Explicit values on the jail refine that attachment. Use \setting{ip\_cidr}
when the prefix must travel with the address, or omit both address keys to use
the network lease allocator. Omit \setting{mac\_address} for a deterministic
address derived from the jail name and bridge. The backend may be
\setting{epair} or \setting{netgraph}; VLAN PVIDs are supported by the epair
backend on a filtering bridge.

Set \setting{shared = true} for host-network inheritance instead of VNET. This
is a weaker isolation mode and is ignored when VNET is enabled.

DNS mode \setting{inherit} copies the host resolver configuration. Mode
\setting{custom} writes the declared nameservers, optional search list, and
domain into the jail.

\section{Health checks}

All checks are command-oriented: exit zero passes and non-zero fails. The
target determines whether the command runs on the host or inside the jail.

\begin{lstlisting}[language=TOML]
[jails.healthcheck]
enabled = true

[[jails.healthcheck.checks]]
name = "http"
command = "fetch -qo /dev/null http://127.0.0.1:8080/health"
target = "jail"
start_period = 60
interval = 30
timeout = 10
retries = 3

[[jails.healthcheck.checks]]
name = "listener"
command = "nc -z 127.0.0.1 8080"
target = "jail"
\end{lstlisting}

The supervisor retains failure history, applies configured recovery behavior,
and can use health as the trigger for a QoS phase transition.

\section{Hooks}

\begin{lstlisting}[language=TOML]
[[jails.hooks]]
phase = "pre_start"
target = "host"
command = "/usr/local/bin/fetch-secrets"
args = ["${jail_name}"]
on_failure = "abort"

[[jails.hooks]]
phase = "pre_stop"
target = "jail"
command = "/usr/sbin/service"
args = ["nginx", "stop"]
on_failure = "continue"
\end{lstlisting}

A hook executes \setting{command} directly with \setting{args} as its argument
vector; there is no shell, so \setting{command = "service nginx stop"} names a
program that does not exist. Use \setting{command = "/bin/sh"} with
\setting{args = ["-c", "..."]} when shell syntax is needed. Arguments expand
\setting{\$\{jail\_name\}}, \setting{\$\{jail\_path\}},
\setting{\$\{jail\_ip\}}, and \setting{\$\{jid\}}.

Supported lifecycle phases include \setting{pre\_create},
\setting{post\_create}, \setting{pre\_start}, \setting{post\_start},
\setting{pre\_stop}, \setting{post\_stop}, \setting{pre\_eva}, and
\setting{post\_eva}. A phase that requires a running jail is rejected when no
JID is available. A \setting{post\_start} hook runs after the jail is up, so
\setting{abort} there only warns: the jail keeps running. Use
\setting{continue} only when failure is genuinely safe to ignore.

\section{Two-phase QoS}

\begin{lstlisting}[language=TOML]
[jails.qos]
transition = "healthy"
max_startup_secs = 120

[jails.qos.startup]
memory = "4g"
cpu = "400"

[jails.qos.steady]
memory = "2g"
cpu = "200"
\end{lstlisting}

\setting{transition} may be \setting{healthy}, \setting{timer}, or
\setting{manual}. Health-driven transition requires the supervisor;
\setting{max\_startup\_secs} provides a deadline fallback. The current phase is
recorded so EVA reconciles against it rather than reverting to startup policy.

\section{Draining and cascade policy}

A drain policy blocks new TCP connections through a jail-specific PF
sub-anchor, then waits up to its timeout for established states to close before
teardown. This is useful for service upgrades behind a load balancer.

\begin{lstlisting}[language=TOML]
[jails.drain]
enabled = true
timeout_secs = 30
\end{lstlisting}

Set \setting{cascade = "strict"}, \setting{"detach"}, or
\setting{"isolate"} on the jail entry to control dependents during
\command{down}:

\begin{description}
  \item[strict] stop dependents first; this is the default;
  \item[detach] stop only the selected target; and
  \item[isolate] leave dependents running but fence them from the stopped target.
\end{description}

\section{Telemetry}

\begin{lstlisting}[language=TOML]
[telemetry]
endpoint = "http://127.0.0.1:4318"
service_name = "blackship"
metrics_interval_secs = 30

[telemetry.headers]
authorization = "Bearer replace-me"
\end{lstlisting}

When compiled with telemetry support, the block enables OTLP/HTTP lifecycle
events and resource gauges. An unavailable collector warns but does not fail a
jail operation. Do not commit credentials in a shared manifest; merge a
protected local override when necessary.

\section{Validation discipline}

Always run both syntax/semantic validation and a lifecycle preview:

\begin{lstlisting}[language=bash]
blackship check
blackship armada config --show
blackship up --all --dry-run
\end{lstlisting}

The first catches manifest and dependency errors, the second reveals merged
Armada configuration, and the third resolves actual operational ordering.
